\documentclass[11pt]{article}
\usepackage{setspace}
\usepackage{geometry}
\geometry{margin=0.85in}
\setstretch{1.0}

\begin{document}

\begin{flushright}
\textbf{Zhongli Wang} \\
+49 160 9541 6015 \\
nkwangzhongli@gmail.com \\
Bielefeld, Germany \\
\end{flushright}

\vspace{0.1cm}

\noindent
\textbf{Julia Mattoscio} \\
Project Lead, ETH Juniors \\
Z\"urich, Switzerland \\
Email: julia.mattoscio@ethjuniors.ch \\

\vspace{0.1cm}

\noindent
\textbf{Application for Data Scientist / Machine Learning Engineer} \\

\vspace{0.05cm}

\noindent Dear Ms Mattoscio,

I have a strong interest in building machine learning systems that move beyond research prototypes and create real value in complex industrial environments. The position advertised through ETH Juniors is especially attractive to me because it combines AI, data science, software development, and operational decision-making. I am currently completing my Ph.D.\ in Economics at Bielefeld University and Universit\`a Cattolica del Sacro Cuore, with planned completion in September 2026. My work sits at the intersection of computational modelling, machine learning, simulation, and applied data science.

Across my doctoral research, I have developed end-to-end workflows in Python for complex systems with heterogeneous behaviour and dynamic responses. This includes data cleaning, feature engineering, reproducible dataset construction, model estimation, validation, large-scale simulation, and robustness checks. My third dissertation paper estimates an agent-based computational model using financial and operational data from AT\&T and T-Mobile, applying the Simulated Method of Moments and neural network-based surrogate modelling. This experience is closely related to industrial data science tasks such as time-series analysis, optimization, simulation-based decision support, and model validation under uncertainty.

I also bring hands-on experience in machine learning, anomaly detection, and GenAI-related workflows. In my co-authored paper \emph{Polytope Fraud Theory}, I worked with supervised learning methods to detect structured irregularities in corporate financial data, strengthening my experience in feature construction, high-dimensional data handling, model evaluation, and interpretation of algorithmic outputs. In \emph{Agentic Workflows for Economic Research}, I worked on structured LLM-assisted workflows for document processing, research automation, code generation, and reproducible data work, with emphasis on auditability, versioning, and human-in-the-loop validation.

The position matches several of my strengths: anomaly detection, time-series and operational analytics, optimization, and Generative/Agentic AI. I am particularly motivated by the opportunity to work closer to production systems. I would be eager to collaborate with software engineers, domain experts, and product teams across the development lifecycle, from problem definition and model development to validation, integration, monitoring, and continuous improvement.

I would be excited to contribute to the industrial technology client by developing robust machine learning models, building reproducible data pipelines, and helping translate AI concepts into operational solutions. My academic training has taught me to work independently on complex problems, communicate assumptions clearly, and collaborate across disciplines. I have very good English skills and strong German skills, which would support communication in an international industrial environment.

Thank you so much for considering my application! I would welcome the opportunity to discuss how my background can contribute to industrial AI and data science projects.

\vspace{0.15cm}

\noindent Sincerely, \\[0.5cm]
\noindent Zhongli Wang

\end{document}
